\documentclass[a4paper,10pt]{article}
\usepackage[a4paper,margin=1.0in]{geometry}
\usepackage{lmodern} % Latin Modern font for pdfLaTeX
% \usepackage{fontspec} % Improved font encoding
\usepackage{xcolor}
\usepackage{enumitem}
\usepackage{titlesec}
\usepackage[
  colorlinks=true,
  linkcolor=highlight,
  citecolor=highlight,
  urlcolor=highlight
]{hyperref}
\usepackage{fontawesome5} % Icons for contact details

% Font and color settings
\renewcommand{\familydefault}{\sfdefault} % Set main font to sans-serif
\definecolor{highlight}{HTML}{4b3069} % Purple highlight color

\pagestyle{empty}

% Commands
% ========
% Commands
% ========

% ========
% \cvevent
% Usage:
% #1: Title
% #2: Date
% #3: Information

\newcommand{\cvevent}[3]{
    \item \textbf{#1} \hfill #2\\
    {#3}
}

% ======
% \paper
% Usage:
% #1: Authors (remember to bold myself!)
% #2: Year
% #3: Title
% #4: Journal reference (e.g. MNRAS, 535(3), 2246-2259)
% #5: DOI (e.g. 10.1093/rasti/rzaf044)


\newcommand{\paper}[5]{
    \item #1 (#2). \textit{#3}. #4.
    \href{https://doi.org/#5}{doi:#5}
}

% Define section format
\titleformat{\section}{\Large\bfseries\color{highlight}}{}{}{}[\titlerule]
\titleformat{\subsection}[runin]{\bfseries}{}{0em}{}[:]

% Customize itemize spacing
\setlist[itemize]{noitemsep, topsep=0pt}

% Header
\def\name{Xander Byrne}
\def\position{Astrophysics PhD Student}
\def\email{
  \href{mailto:xbyrne@ast.cam.ac.uk}{xbyrne@ast.cam.ac.uk}
}
\def\website{
  \href{https://xbyrne.github.io}{xbyrne.github.io}
}
\def\address{
  Institute of Astronomy,
  Cambridge, UK
}
\def\orcid{
  \href{https://orcid.org/0000-0001-9488-238X}{0000-0001-9488-238X}
}

\begin{document}

% Top header
\noindent
\begin{minipage}[t]{0.73\textwidth}
  \vspace{1.2ex} % Adjust spacing as necessary
  {\fontsize{50}{55}\selectfont \textbf{\name}}\\[0.3em]
  {\fontsize{20.5}{30}\textit{\position}}
\end{minipage}
{\color{highlight}\vrule width 1.5pt}
\begin{minipage}[t]{0.27\textwidth}
  \raggedleft
  {\fontsize{11}{12}\selectfont
    \vspace{.3em}
    \address \faMapMarker* \\[.3em]
    \email \faEnvelope[regular] \\[.1em]
    \website \faGithub \\[.3em]
    \orcid \faOrcid
  }
\end{minipage}

% Section: Research Interests
\section*{Research Interests}
$\bullet$ Exoplanet composition and habitability
\hspace{1.5em}
$\bullet$ Polluted white dwarfs
\hspace{1.5em}
$\bullet$ Data science \& machine learning

\section*{Education}

\textbf{Ph.D. in Astronomy}, Institute of Astronomy, University of Cambridge \hfill 2023---\\
Dissertation: \textit{Clues to exoplanet composition from stellar spectra}.\\
Supervisor: Dr Amy Bonsor\\

\noindent \textbf{M.Sci in Natural Sciences (Astrophysics)}, University of Cambridge \hfill 2019--23 \\
Graduated 1st class with distinction.

\section*{Publications}

\begin{itemize}[leftmargin=8pt]

    \paper{\textbf{X.~Byrne} \& A.~Bonsor}{2026, subm.}{
      The true masses of radial-velocity exoplanets constrained by stability.
    }{RASTI}{TBD}

    \paper{\textbf{X.~Byrne}, C.~M.~Guimond, A.~Bonsor, et al.}{2026}{
      The Barnard's Star Planetary System: stability, composition, and evolution of four sub-Earth exoplanets.
    }{MNRAS}{10.1093/mnras/stag1207}

    \paper{\textbf{X.~Byrne}, A.~Bonsor, L.~K.~Rogers et al.}{2025}{
      Finding rare classes in large datasets: the case of polluted white dwarfs from \textit{Gaia} XP spectra.
    }{RASTI, 4}{10.1093/rasti/rzaf044}

    \paper{\textbf{X.~Byrne}, A.~Bonsor, L.~K.~Rogers et al.}{2024}{
      Semi-supervised spectral classification of DESI white dwarfs by dimensionality reduction.
    }{MNRAS, 535(3), 2246-2259}{10.1093/mnras/stae2478}

    \paper{\textbf{X.~Byrne}, R.~A.~Meyer, E.~P.~Farina, et al.}{2024}{
      Quasar Island - three new $z\sim6$ quasars, including a lensed candidate, identified with contrastive learning.
    }{MNRAS, 530(1), 870-880}{10.1093/mnras/stae902}

    \paper{\textbf{X.~Byrne}, O.~Shorttle, S.~Jordan et al.}{2024}{
      Atmospheres as a window to rocky exoplanet surfaces.
    }{MNRAS, 527(4), 10748-10759}{10.1093/mnras/stad3914}
\end{itemize}

\section*{Talks}

\subsection{Invited Talks}

\begin{itemize}
    \vspace{1ex}
    \cvevent{White Dwarf Meeting -- University of Warwick}{Feb 2026}{
      \textit{Big-Data Methods for Polluted WDs -- Past, Present and Future}
    }

    \cvevent{ExoCoffee -- MPIA, Heidelberg}{Feb 2024}{
      \textit{Atmospheres as a Window to Rocky Planet Surfaces}
    }

    \cvevent{Data Intensive Science Seminar Series -- Maxwell Centre, Cambridge}{Feb 2024}{
      \textit{Contrastive Learning and High-Redshift Quasars}
    }

\end{itemize}

\subsection{Contributed Talks}

\begin{itemize}
    \vspace{1ex}

    \cvevent{EuroWD -- ISTA, Vienna}{Aug 2026}{
      \textit{Big-data methods for polluted white dwarfs: past, present, and future}
    }

    \cvevent{Origins 2026 -- Paris}{Jul 2026}{
      \textit{The Barnard's Star Planetary System: interior structure and water capacity of four sub-Earth planets}
    }

    \cvevent{UK White Dwarf Science Meeting -- Queens' College, Cambridge}
    {Mar 2026}{
      \textit{Calcium Needles in Hydrogen Haystacks -- Identifying polluted white dwarfs with \textit{Gaia}}
    }

    \cvevent{Wednesday Seminar Series -- Institute of Astronomy, Cambridge}{Oct 2024}{
      \textit{Spectral Classification of White Dwarfs by Dimensionality Reduction}
    }

    \cvevent{Astro Data Science Discussion Group -- KICC, Cambridge}{Jan 2024}{
      \textit{Contrastive Learning and High-Redshift Quasars}
    }

    \cvevent{Galaxy Coffee -- MPIA, Heidelberg}{Sep 2022}{
      \textit{Searching for Gravitationally-Lensed High-Redshift Quasars with Unsupervised Machine Learning}
    }
\end{itemize}

\section*{Awarded Telescope Time}
% \subsection*{As PI}

\subsection*{As Co-I}

\begin{itemize}
  \item \textit{White dwarfs as a means to probe core formation in exoplanetary systems}; PI Rogers\\
    65h VLT/X-Shooter spectroscopy (115.28F4)

  \item \textit{Characterising a $z\sim6$ gravitationally-lensed BAL quasar candidate}; PI Meyer\\
    2.9h VLT/X-Shooter spectroscopy (113.26CY)

  \item \textit{Measuring the mass of the SMBHs powering two lensed $z\sim6$ quasars}; PI Meyer\\
    4.23h Gemini (DDT) FLAMINGOS-2 spectroscopy (DT-2022B-029)
\end{itemize}
\section*{Professional Service}

\subsection{Conferences \& Seminars}

\begin{itemize}
    \cvevent{Organiser, UK White Dwarf Science Meeting}{Mar 2026}{
      Coordinating invitations, payments, scheduling, and talks for a national one-day meeting.
    }

    \cvevent{Organiser, Wednesday Seminar Series}{2024---}{
      A long-running seminar series given by PhD students and postdocs at the Institute of Astronomy.
    }

    \cvevent{Organiser, PhD/Postdoc Jamboree}{2024---}{
      A social `introduction' event for the Institute's PhDs and postdocs.
    }

    \cvevent{Tech Assistant, \textit{IoA50} conference}{Jul 2024}{
      Assisting speakers at a large conference with presentations, mics etc.
    }
\end{itemize}

\subsection{Other}

\begin{itemize}
    \cvevent{Co-chair, LGBTQ+ Focus Group, Institute of Astronomy}{2025--26}{
      Organising departmental events.
    }

    \cvevent{Organiser, Undergraduate Journal Club}{2023--24}{
      A journal club for undergraduates, also offering career advice.
    }
\end{itemize}

\section*{Grants, Awards \& Prizes}
\begin{itemize}
  \item \textbf{IoA Community Service Award} \hfill Institute of Astronomy\\
    For `exemplary contributions to Public Outreach and the Wednesday Seminar Committee'.
  \item \textbf{MASt/Part III Bursary} (£3500) \hfill Institute of Astronomy\\
    Funding the conversion of my Masters' thesis to a publication over the course of a summer internship.
  \item \textbf{MPIA Summer Internship} (€3750) \hfill Max-Planck-Institut f\"{u}r Astronomie\\
    Funding internship: \textit{Searching for Lensed High-Redshift Quasars using Unsupervised Machine Learning}.
  \item \textbf{Winifred Georgia Holgate-Pollard Memorial Prize} (£1650)
    \hfill University of Cambridge\\
    Awarded for the best performance in the Natural Sciences (Astrophysics) examinations.
  \item \textbf{John Spencer Wilson Prize} (£500) \hfill St Catharine's College, Cambridge\\
    Awarded in recognition of exceptional performance in the Natural Sciences examinations.
\end{itemize}

\section*{Teaching}

\noindent \textbf{Supervisor}, Part II Astrophysics, University of Cambridge \hfill 2023---
\begin{itemize}
  \item Conducted small-group teaching, marking solutions
    \begin{itemize}
      \item \textit{Relativity} (GR): 18 students
      \item \textit{Stellar Dynamics \& Structure of Galaxies} (orbits \& potentials): 8 students
    \end{itemize}
  \item Produced \href{https://xbyrne.github.io/notes}{popular revision notes} ($\sim2500$ hits per year)
\end{itemize}

\section*{Miscellaneous Experience}

\subsection{Technical Experience}

\begin{itemize}
  \item \textbf{Programming Languages:} Python, C++, MATLAB

  \item \textbf{Astronomy Software:} EsoRex, EsoReflex, DS9, TOPCAT, PyTorch, TensorFlow

  \item \textbf{Development:} Primary developer and maintainer of the open-source \href{https://github.com/xbyrne/cinemas}{\texttt{cinemas} package}

  \item \textbf{Datasets:} Gaia, DESI, SDSS, GALAH
\end{itemize}

\subsection{Industry Experience}

\begin{itemize}
    \cvevent{Internship at \href{https://www.featurespace.com}{Featurespace}}{Feb--Aug 2025}{
      6-month placement in the R\&D division of a tech company.
      Implemented and tested state-of-the-art innovations in AI research into the company's extensive codebase.
    }
\end{itemize}

\section*{Outreach}

\begin{itemize}
    \cvevent{IoA+KICC Open Day}{2024---}{
      Organising and running activities, talks, and demonstrations for the general public.
    }

    \cvevent{\textit{The Naked Scientists}, BBC Radio 5 Live}{Mar 2023---}{
      National radio show and podcast. Numerous appearances both live and in pre-recorded segments
    }

    \cvevent{Public Open Evenings, Institute of Astronomy}{2022---}{
      Visits for the general public, involving a public lecture (which I gave in \href{https://www.youtube.com/watch?v=sq6PVKXRj4s}{2025} and \href{https://www.youtube.com/watch?v=fqoo_07Zfsg}{2026}) and stargazing.
    }

    \cvevent{Asdan Summer School, St Catharine's College, Cambridge}{Aug 2025}{
      Outreach talk: \textit{Worlds Anew: Exoplanetary Astronomy}
    }

    \cvevent{Cambridge Hands-On Science (CHaOS) event: \textit{Crash Bang Squelch}}{Mar 2024}{
      Outreach talk: \textit{How to find an Alien Planet}
    }

    \cvevent{Birmingham Schools Science Network}{Sep 2023}{
      Outreach talk: \href{https://www.youtube.com/watch?v=J-r3r0u_INQ}{\textit{Worlds Anew: Quasars, Exoplanets, and the Search for Life}}.
    }

    \cvevent{HiSPARC Project}{2018--22}{
      Teaching sessions and sparking research projects for high-school physics students, particularly girls
    }

    \cvevent{\textit{Varsity Sci 2021}}{Oct 2021}{
      Outreach talk: \href{https://www.youtube.com/live/_3y9bzUpaeY?si=sBWpn4ZTFP_lCF9S&t=3638}{\textit{The Messier Catalogue: 110 objects in 10 minutes}}.
    }

    \cvevent{Cambridge University Physics Society}{Aug 2020}{
      Outreach talk: \href{https://youtu.be/ovptNIUAtto?si=m9zADxlmFnNeBend}{\textit{Negative mass: matter that's lighter than light}}.
    }
\end{itemize}

\section*{Press Coverage}

\begin{itemize}
  \item \textbf{Institute of Astronomy}:
    \href{https://www.ast.cam.ac.uk/news/astronomers-find-nearby-planets-be-small-strange-and-utterly-uninhabitable}{
      \textit{Astronomers find nearby planets to be small, strange, and}
    } \hfill Jun 2026\\
    \href{https://www.ast.cam.ac.uk/news/astronomers-find-nearby-planets-be-small-strange-and-utterly-uninhabitable}{
      \textit{utterly uninhabitable}
    }

  \item \textbf{NOIRLab Stories Blog}:
    \href{https://noirlab.edu/public/blog/cosmic-gems/}{
      \textit{AI Aids in Search for Cosmic Gems in Ocean of Data}
    }
    \hfill Jul 2024
\end{itemize}

\end{document}